\section{Conclusion}
\label{sec:conclusion}

The Portuguese minimum wage has had two distinct lives. For its first four
decades it lost ground: inflation eroded its real value faster than statutory
increases restored it, and it fell far behind an economy whose output per worker
more than doubled. Measured against productivity it reached barely a third of
its introductory value. Only in the last decade has that reversed, with the
policy residual turning persistently positive from 2015 and the real floor
regaining its 1974 level around 2017.

That is a descriptive account, and we have been careful to keep it one. The
residual measures what the wage floor did relative to a productivity-and-
inflation benchmark; it does not measure what the wage floor caused. We have
also shown, rather than asserted, why the natural causal design is unavailable:
the Azorean supplement is proportional and contributes no independent timing,
Madeira is the sole source of genuine regional divergence, and with the small
number of Portuguese regions the conventional inference that would make such a
design look conclusive rejects a true null far too often. On the actual data,
estimates significant at the one-tenth of one per cent level under clustered
standard errors survive none of a wild cluster bootstrap.

We think the honest conclusion is that Portugal is a poor setting for
identifying minimum-wage price pass-through from regional variation, and that
saying so is more useful than a weakly identified estimate. The country's
statutory history is, however, unusually well documented in primary law, and the
long-run record it supports is informative in its own right: a wage floor whose
real generosity has moved by a factor of two over fifty years, in both
directions, is a substantial object of study.

The series, the regional schedules, the price panel and the pipeline that builds
them are released with this paper. Registering the Madeira acts that remain
outside it, and obtaining minimum-wage coverage jointly by region and activity,
would be the two most valuable next steps for anyone wishing to revisit the
causal question.
